% Unit 1 self-study -- "I do / You do" ladder for CED 1.1-1.8, four YOUR
% TURN questions per ladder. Ten ladders, forty questions.
%
% Values are chosen DISJOINT from u01-notes, the u01 worksheets and
% u01-frq: the FRQ set uses boron isotopes, an aluminium PES spectrum, an
% acetic-acid combustion and Ca3N2/CaO-KCl, so this document uses copper
% and lithium isotopes, a phosphorus PES spectrum, percent-composition
% routes to formulas, and Mg3N2/MgO-NaF instead. Nothing is repeated.
\documentclass{apchem-worksheet}

\apcunit{1}
\apcunittitle{Atomic Structure and Properties}
\apcdoctitle{Self-Study \textbullet\ CED 1.1--1.8, I~do / You~do}
\apcsubtitle{Ten ladders \textbullet\ four YOUR TURN questions each \textbullet\ work alone, check after all four}

\begin{document}
\apctitleblock

\begin{apcnote}[How to use these notes --- read this first]
Each skill is a \textbf{ladder}: a fully worked example in the
solid-framed box, then \textbf{four} YOUR TURN questions in the dashed box
--- same skill, new numbers, no help.

\begin{enumerate}[leftmargin=*,itemsep=0.15em]
  \item \textbf{Work all four} before checking anything.
  \item Check against the gray \emph{check:} line. All four right first
        time earns the tick on the tracker (last page).
  \item Any misses: re-read the worked example, then redo the missed ones
        from a blank page --- not by reading the answer again.
\end{enumerate}

\textbf{Two habits that prevent most lost marks in this unit.} Always
convert to \textbf{moles} before doing anything else with a mass ---
almost every Unit 1 error is a comparison made in grams that should have
been made in moles. And when you explain a periodic trend, name
\textbf{nuclear charge, distance and shielding}. An answer that says an
atom ``wants a full shell'' earns nothing on the AP exam.
\end{apcnote}

% =====================================================================
\section*{\sffamily Ladder 1 \textbullet\ Moles, mass and particles}
\ced{1.1}

The mole is a counting unit: one mole is
$6.022\times10^{23}$ particles. The molar mass in
\si{\gram\per\mole} is the bridge between the mass you can weigh and the
number of particles you cannot.

\begin{workedexample}[I do: mass to moles to atoms]
How many moles, and how many atoms, are in \SI{12.5}{\gram} of calcium?

\textbf{Mass to moles} --- divide by the molar mass:
\[ n = \frac{12.5}{40.08} = \mathbf{0.312~mol} \]

\textbf{Moles to atoms} --- multiply by Avogadro's number:
\[ 0.3119 \times 6.022\times10^{23} = \mathbf{1.88\times10^{23}}
   \text{ atoms} \]

\textbf{Sense check:} \SI{12.5}{\gram} is well under one molar mass
(\SI{40.08}{\gram}), so the answer must be under one mole. It is.
\end{workedexample}

\begin{yourturn}[YOUR TURN 1 --- four questions]
\begin{enumerate}[leftmargin=*,itemsep=0.5em]
  \item Moles in \SI{8.40}{\gram} of iron: \workspace[1.4cm]
  \item Atoms in \SI{2.25}{\gram} of magnesium: \workspace[1.6cm]
  \item Mass of \SI{3.01e23}{} atoms of sodium: \workspace[1.6cm]
  \item Without calculating: \SI{10}{\gram} of helium or \SI{10}{\gram} of
        iron --- which contains more atoms, and why? \workspace[1.6cm]
\end{enumerate}
\selfcheck{(a) \num{0.150}~mol \quad (b) \num{5.57e22} atoms \quad
(c) \SI{11.5}{\gram} \quad (d) helium --- far smaller molar mass, so the
same mass is many more moles}
\keyonly{\begin{apcanswer}(a) $8.40/55.85 = \mathbf{0.150}$~mol.
(b) $2.25/24.31 = 0.0926$~mol; $\times 6.022\times10^{23} =
\mathbf{5.57\times10^{22}}$ atoms.
(c) $3.01\times10^{23}/6.022\times10^{23} = 0.500$~mol;
$\times 22.99 = \mathbf{11.5}$~g.
(d) \textbf{Helium.} Molar mass \num{4.003} against \num{55.85}, so
\SI{10}{\gram} of helium is about 2.5~mol while \SI{10}{\gram} of iron is
about 0.18~mol --- roughly fourteen times as many atoms. Equal
\emph{masses} are not equal counts.\end{apcanswer}}
\end{yourturn}

% =====================================================================
\section*{\sffamily Ladder 2 \textbullet\ Molar mass of a compound}
\ced{1.1}

A subscript outside a bracket multiplies \emph{everything} inside it.
That single rule is where most molar-mass errors come from.

\begin{workedexample}[I do: a compound with brackets]
Find the molar mass of \ce{Ca(NO3)2}.

\textbf{Read the formula first.} There is one calcium, and \emph{two}
nitrate groups. Each nitrate is one N and three O, so the compound holds
2 N and 6 O --- not 1 N and 3 O.

\[ 40.08 + 2\,[\,14.01 + 3(16.00)\,] = 40.08 + 2(62.01)
   = \mathbf{164.1~g/mol} \]

\textbf{The trap:} forgetting to distribute the 2 gives \num{102.1}, and
every later answer built on it is wrong.
\end{workedexample}

\begin{yourturn}[YOUR TURN 2 --- four questions]
\begin{enumerate}[leftmargin=*,itemsep=0.5em]
  \item Molar mass of \ce{MgCl2}: \workspace[1.4cm]
  \item Molar mass of \ce{K2CO3}: \workspace[1.4cm]
  \item Molar mass of \ce{Al2(SO4)3}: \workspace[1.8cm]
  \item How many oxygen atoms are in one formula unit of
        \ce{Al2(SO4)3}? \workspace[1.4cm]
\end{enumerate}
\selfcheck{(a) \num{95.21} \quad (b) \num{138.2} \quad (c) \num{342.1}
\quad (d) 12}
\keyonly{\begin{apcanswer}(a) $24.31 + 2(35.45) = \mathbf{95.21}$~g/mol.
(b) $2(39.10) + 12.01 + 3(16.00) = \mathbf{138.2}$~g/mol.
(c) $2(26.98) + 3\,[\,32.06 + 4(16.00)\,] = 53.96 + 3(96.06) =
\mathbf{342.1}$~g/mol.
(d) Three sulfate groups, four oxygens each: $3 \times 4 = \mathbf{12}$
oxygen atoms.\end{apcanswer}}
\end{yourturn}

% =====================================================================
\section*{\sffamily Ladder 3 \textbullet\ Mass spectra and average atomic mass}
\ced{1.2}

A mass spectrum plots mass on the horizontal axis and abundance on the
vertical. Peak \emph{position} gives isotope mass; peak \emph{height}
gives how common that isotope is. The tabulated atomic mass is the
abundance-weighted average of the two.

\begin{workedexample}[I do: two isotopes to one atomic mass]
Copper has two isotopes: \ce{^{63}Cu},
\SI{62.930}{\atomicmassunit}, \SI{69.17}{\percent}; and \ce{^{65}Cu},
\SI{64.928}{\atomicmassunit}, \SI{30.83}{\percent}. Find the average
atomic mass.

\textbf{Weight each mass by its fractional abundance and add:}
\[ (62.930)(0.6917) + (64.928)(0.3083) = 43.53 + 20.02 =
   \mathbf{63.55~u} \]

\textbf{Sense check:} the answer must land between the two isotope masses,
and closer to the \emph{more abundant} one. \num{63.55} sits nearer
\num{62.930}, as it must with \SI{69}{\percent} of atoms being the lighter
isotope.

\textbf{The trap:} averaging \num{62.930} and \num{64.928} without
weighting gives \num{63.93} --- wrong, because it assumes equal
abundances.
\end{workedexample}

\begin{yourturn}[YOUR TURN 3 --- four questions]
\begin{enumerate}[leftmargin=*,itemsep=0.5em]
  \item Lithium: \SI{6.015}{\atomicmassunit} at \SI{7.59}{\percent},
        \SI{7.016}{\atomicmassunit} at \SI{92.41}{\percent}. Average?
        \workspace[1.6cm]
  \item Chlorine: \SI{34.969}{\atomicmassunit} at \SI{75.78}{\percent},
        \SI{36.966}{\atomicmassunit} at \SI{24.22}{\percent}. Average?
        \workspace[1.6cm]
  \item An element has isotopes at 20 u and 22 u and an average atomic
        mass of \num{20.2}. Which peak is taller, and roughly by how much?
        \workspace[1.6cm]
  \item Explain why no single atom of chlorine has a mass of
        \SI{35.45}{\atomicmassunit}. \workspace[1.8cm]
\end{enumerate}
\selfcheck{(a) \num{6.94} \quad (b) \num{35.45} \quad (c) the 20 u peak,
by roughly 9:1 \quad (d) it is an average over isotopes, not a real
particle mass}
\keyonly{\begin{apcanswer}(a) $(6.015)(0.0759) + (7.016)(0.9241) =
\mathbf{6.94}$~u.
(b) $(34.969)(0.7578) + (36.966)(0.2422) = \mathbf{35.45}$~u.
(c) The \textbf{20 u} peak is much taller. The average sits only
\num{0.2} above 20 but \num{1.8} below 22, so the light isotope must
dominate --- about \SI{90}{\percent} against \SI{10}{\percent}, a ratio
near $9\!:\!1$.
(d) Chlorine atoms are either \ce{^{35}Cl} or \ce{^{37}Cl}. The tabulated
\num{35.45} is the \textbf{abundance-weighted average} of those two, a
property of a bulk sample rather than of any individual
atom.\end{apcanswer}}
\end{yourturn}

% =====================================================================
\section*{\sffamily Ladder 4 \textbullet\ Percent composition}
\ced{1.3}

Percent composition asks what fraction of a compound's mass each element
contributes. Work from \emph{one mole} of the compound.

\begin{workedexample}[I do: percent composition of a nitrate]
Find the percent by mass of each element in \ce{Ca(NO3)2}
($M = \SI{164.1}{\gram\per\mole}$).

\textbf{Mass of each element in one mole:}
Ca $= 40.08$; \quad N $= 2(14.01) = 28.02$; \quad
O $= 6(16.00) = 96.00$.

\[ \%\text{Ca} = \frac{40.08}{164.1}\times100 = \mathbf{24.4\%} \qquad
   \%\text{N} = \frac{28.02}{164.1}\times100 = \mathbf{17.1\%} \]
\[ \%\text{O} = \frac{96.00}{164.1}\times100 = \mathbf{58.5\%} \]

\textbf{Sense check:} the three must sum to \SI{100}{\percent}.
$24.4 + 17.1 + 58.5 = 100.0$. \checkmark
\end{workedexample}

\begin{yourturn}[YOUR TURN 4 --- four questions]
\begin{enumerate}[leftmargin=*,itemsep=0.5em]
  \item Percent chlorine in \ce{MgCl2}: \workspace[1.6cm]
  \item Percent potassium in \ce{K2CO3}: \workspace[1.6cm]
  \item Percent oxygen in \ce{K2CO3}: \workspace[1.6cm]
  \item Which has the higher percent by mass of oxygen, \ce{H2O} or
        \ce{H2O2}? Justify without a full calculation.
        \workspace[1.8cm]
\end{enumerate}
\selfcheck{(a) \SI{74.5}{\percent} \quad (b) \SI{56.6}{\percent} \quad
(c) \SI{34.7}{\percent} \quad (d) \ce{H2O2} --- twice the oxygen for only
the same hydrogen}
\keyonly{\begin{apcanswer}(a) $2(35.45)/95.21 \times 100 =
\mathbf{74.5}\%$.
(b) $2(39.10)/138.2 \times 100 = \mathbf{56.6}\%$.
(c) $3(16.00)/138.2 \times 100 = \mathbf{34.7}\%$.
(d) \textbf{\ce{H2O2}}. Going from \ce{H2O} to \ce{H2O2} doubles the
oxygen mass while the hydrogen mass is unchanged, so oxygen must take a
larger share. Confirming: \ce{H2O} is \SI{88.8}{\percent} O and
\ce{H2O2} is \SI{94.1}{\percent} O.\end{apcanswer}}
\end{yourturn}

% =====================================================================
\section*{\sffamily Ladder 5 \textbullet\ Empirical and molecular formulas}
\ced{1.3}

The empirical formula is the simplest whole-number \emph{ratio} of atoms.
The molecular formula is the actual count. Percent composition gives you
the first; only the molar mass gives you the second.

\begin{workedexample}[I do: percentages to a molecular formula]
A compound is \SI{40.00}{\percent} C, \SI{6.71}{\percent} H and
\SI{53.29}{\percent} O by mass, with
$M = \SI{180.2}{\gram\per\mole}$.

\textbf{Step 1 --- assume \SI{100.0}{\gram},} so the percentages become
masses in grams.

\textbf{Step 2 --- convert each to moles:}
\[ \text{C: } \frac{40.00}{12.011} = 3.331 \qquad
   \text{H: } \frac{6.71}{1.008} = 6.657 \qquad
   \text{O: } \frac{53.29}{16.00} = 3.331 \]

\textbf{Step 3 --- divide by the smallest} (\num{3.331}):
C~1.00, H~2.00, O~1.00, so the empirical formula is \textbf{\ce{CH2O}}.

\textbf{Step 4 --- scale to the molar mass.} Empirical mass
$= 12.01 + 2(1.008) + 16.00 = 30.03$, and
$180.2/30.03 = 2$\ldots{} in fact $\mathbf{6.00}$.
\[ \ce{CH2O} \times 6 = \mathbf{\ce{C6H12O6}} \]

\textbf{The trap:} stopping at \ce{CH2O}. That ratio is shared by
formaldehyde, acetic acid and glucose --- the molar mass is what separates
them.
\end{workedexample}

\begin{yourturn}[YOUR TURN 5 --- four questions]
\begin{enumerate}[leftmargin=*,itemsep=0.5em]
  \item \SI{92.26}{\percent} C, \SI{7.74}{\percent} H,
        $M = \SI{78.11}{\gram\per\mole}$. Empirical and molecular
        formulas? \workspace[2.2cm]
  \item \SI{5.93}{\percent} H, \SI{94.07}{\percent} O,
        $M = \SI{34.01}{\gram\per\mole}$. Empirical and molecular?
        \workspace[2.2cm]
  \item \SI{69.94}{\percent} Fe, \SI{30.06}{\percent} O. Empirical
        formula? \workspace[2.0cm]
  \item Why can two different compounds share an empirical formula but
        never a molecular formula \emph{and} the same structure?
        \workspace[1.6cm]
\end{enumerate}
\selfcheck{(a) \ce{CH}, \ce{C6H6} \quad (b) \ce{HO}, \ce{H2O2} \quad
(c) \ce{Fe2O3} \quad (d) the empirical formula fixes only the ratio}
\keyonly{\begin{apcanswer}(a) C: $92.26/12.011 = 7.681$; H:
$7.74/1.008 = 7.679$. Ratio $1:1$, so \textbf{\ce{CH}} (mass
\num{13.02}). $78.11/13.02 = 6$, giving \textbf{\ce{C6H6}}, benzene.
(b) H: $5.93/1.008 = 5.883$; O: $94.07/16.00 = 5.879$. Ratio $1:1$, so
\textbf{\ce{HO}} (mass \num{17.01}). $34.01/17.01 = 2$, giving
\textbf{\ce{H2O2}}.
(c) Fe: $69.94/55.85 = 1.252$; O: $30.06/16.00 = 1.879$. Dividing by the
smaller: Fe~1.00, O~1.50. Multiply both by 2 to clear the half:
\textbf{\ce{Fe2O3}}.
(d) The empirical formula fixes only the \emph{ratio} of atoms, so any
number of compounds can share it. The molecular formula fixes the actual
count --- and two compounds can still share that (isomers) while differing
in how the atoms are connected.\end{apcanswer}}
\end{yourturn}

% =====================================================================
\section*{\sffamily Ladder 6 \textbullet\ Composition of mixtures}
\ced{1.4}

A mixture has no formula. Its composition is whatever was put in, so it is
described by percent by mass --- and its components keep their own
properties, which is how they can be separated.

\begin{workedexample}[I do: separating and weighing]
A \SI{8.00}{\gram} mixture of \ce{KCl} and sand is stirred with water,
filtered, and the filtrate evaporated to dryness, leaving
\SI{4.60}{\gram} of \ce{KCl}. Find the percent by mass of \ce{KCl}, and
explain why the method works.

\[ \frac{4.60}{8.00}\times100 = \mathbf{57.5\%~KCl} \]

\textbf{Why it works:} \ce{KCl} is ionic and dissolves as hydrated
\ce{K+} and \ce{Cl-} ions; sand is a covalent network solid whose bonds
water cannot break. The difference in \textbf{solubility} is the property
being exploited, and filtration separates the two.
\end{workedexample}

\begin{yourturn}[YOUR TURN 6 --- four questions]
\begin{enumerate}[leftmargin=*,itemsep=0.5em]
  \item A \SI{12.0}{\gram} mixture yields \SI{7.20}{\gram} of soluble
        salt. Percent salt? \workspace[1.4cm]
  \item A mixture is \SI{2.00}{\gram} \ce{NaCl} and \SI{3.00}{\gram}
        \ce{KCl}. What is the percent by mass of \emph{sodium} in the
        whole mixture? \workspace[2.0cm]
  \item Name a technique that separates two miscible liquids, and the
        property it exploits. \workspace[1.6cm]
  \item Why can a mixture's composition vary while a compound's cannot?
        \workspace[1.6cm]
\end{enumerate}
\selfcheck{(a) \SI{60.0}{\percent} \quad (b) \SI{15.7}{\percent} \quad
(c) distillation; difference in boiling point / volatility \quad
(d) mixtures are not chemically combined}
\keyonly{\begin{apcanswer}(a) $7.20/12.0 \times 100 = \mathbf{60.0}\%$.
(b) Sodium in the \ce{NaCl}: $2.00 \times (22.99/58.44) =
\SI{0.787}{\gram}$. As a fraction of the whole mixture:
$0.787/5.00 \times 100 = \mathbf{15.7}\%$. Note the denominator is the
\emph{total} mixture mass, not the \ce{NaCl} mass.
(c) \textbf{Distillation}, exploiting the difference in
\textbf{volatility} (boiling point): the more volatile liquid vaporizes
first and is condensed separately.
(d) A mixture's components are only physically combined in whatever
proportion was used, so the proportion can be anything. A compound's atoms
are chemically bonded in a fixed whole-number ratio, so its composition is
fixed by its formula.\end{apcanswer}}
\end{yourturn}

% =====================================================================
\section*{\sffamily Ladder 7 \textbullet\ Electron configurations}
\ced{1.5}

Fill lowest-energy orbitals first: \ce{1s}, \ce{2s}, \ce{2p}, \ce{3s},
\ce{3p}, \ce{4s}, \ce{3d}, \ce{4p}. Capacities are $s$~2, $p$~6, $d$~10.
The superscripts must sum to the electron count.

\begin{workedexample}[I do: an atom and an ion]
Write the configuration of phosphorus ($Z = 15$) and of the sulfide ion,
\ce{S^2-}.

\textbf{Phosphorus, 15 electrons:}
\[ \ce{1s^2 2s^2 2p^6 3s^2 3p^3} \qquad
   (2+2+6+2+3 = 15 \text{ \checkmark}) \]

\textbf{Sulfide.} Neutral sulfur has 16 electrons; the $2-$ charge means
it has \emph{gained} two, giving 18:
\[ \ce{1s^2 2s^2 2p^6 3s^2 3p^6} \qquad (=18 \text{ \checkmark}) \]

\textbf{The trap:} writing the ion's configuration as though it were the
atom. Always adjust the electron count for the charge first --- negative
means more electrons, positive means fewer.
\end{workedexample}

\begin{yourturn}[YOUR TURN 7 --- four questions]
\begin{enumerate}[leftmargin=*,itemsep=0.5em]
  \item Configuration of calcium ($Z = 20$): \workspace[1.6cm]
  \item Configuration of \ce{Al^3+}: \workspace[1.6cm]
  \item Configuration of iron ($Z = 26$): \workspace[1.8cm]
  \item Explain why \ce{Al^3+} and \ce{Mg^2+} have identical
        configurations but different sizes. \workspace[1.8cm]
\end{enumerate}
\selfcheck{(a) \ce{1s^2 2s^2 2p^6 3s^2 3p^6 4s^2} \quad
(b) \ce{1s^2 2s^2 2p^6} \quad
(c) \ce{1s^2 2s^2 2p^6 3s^2 3p^6 4s^2 3d^6} \quad
(d) same electrons, different nuclear charge}
\keyonly{\begin{apcanswer}(a) \ce{1s^2 2s^2 2p^6 3s^2 3p^6 4s^2}
($= 20$).
(b) Aluminium has 13 electrons; losing three leaves 10:
\ce{1s^2 2s^2 2p^6}.
(c) \ce{1s^2 2s^2 2p^6 3s^2 3p^6 4s^2 3d^6} ($= 26$). Note \ce{4s} fills
before \ce{3d}.
(d) Both are \textbf{isoelectronic} --- ten electrons in the same
arrangement --- so shielding is essentially identical. But aluminium has
13 protons against magnesium's 12. The greater nuclear charge pulls the
same electron cloud in more tightly, so \ce{Al^3+} is the
\textbf{smaller} ion.\end{apcanswer}}
\end{yourturn}

% =====================================================================
\section*{\sffamily Ladder 8 \textbullet\ Reading a photoelectron spectrum}
\ced{1.6}

A PES spectrum has one peak per occupied \emph{subshell}. Peak
\textbf{height} gives the number of electrons in that subshell; peak
\textbf{position} gives their binding energy. Summing the heights gives
the atomic number --- which identifies the element.

\begin{workedexample}[I do: peaks to an element]
A spectrum shows five peaks of relative height $2$, $2$, $6$, $2$ and $3$,
running from highest binding energy to lowest. Identify the element and
explain which peak is highest in binding energy.

\textbf{Sum the heights:} $2+2+6+2+3 = 15$ electrons, so $Z = 15$ ---
\textbf{phosphorus}, \ce{1s^2 2s^2 2p^6 3s^2 3p^3}.

\textbf{Highest binding energy: the first \ce{1s} peak.} Those electrons
lie closest to the nucleus and are shielded by no others, so they feel the
largest effective nuclear charge over the shortest distance. The Coulombic
attraction is strongest, so the most energy is needed to remove one.

\textbf{Lowest binding energy: the \ce{3p} peak of height 3.} Outermost
and most shielded, so held least tightly --- these are the electrons
phosphorus loses or shares in chemistry.
\end{workedexample}

\begin{yourturn}[YOUR TURN 8 --- four questions]
\begin{enumerate}[leftmargin=*,itemsep=0.5em]
  \item Peaks of height $2, 2, 6, 1$. Identify the element.
        \workspace[1.6cm]
  \item Peaks of height $2, 2, 6, 2, 4$. Identify the element.
        \workspace[1.6cm]
  \item An element's spectrum has four peaks, the last of height 2.
        Identify it and give its configuration. \workspace[1.8cm]
  \item Two elements both show a \ce{1s} peak. Explain why the peak sits
        at a \emph{higher} binding energy for the element of larger $Z$.
        \workspace[1.8cm]
\end{enumerate}
\selfcheck{(a) sodium \quad (b) sulfur \quad (c) magnesium,
\ce{1s^2 2s^2 2p^6 3s^2} \quad (d) more protons pull the \ce{1s}
electrons harder}
\keyonly{\begin{apcanswer}(a) $2+2+6+1 = 11$, so $Z = 11$:
\textbf{sodium}.
(b) $2+2+6+2+4 = 16$: \textbf{sulfur}.
(c) $2+2+6+2 = 12$: \textbf{magnesium}, \ce{1s^2 2s^2 2p^6 3s^2}.
(d) The \ce{1s} electrons are essentially unshielded in both, so the
attraction they feel is set almost entirely by the number of protons. More
protons means a larger effective nuclear charge at the same tiny distance,
hence a stronger Coulombic attraction and a higher binding
energy.\end{apcanswer}}
\end{yourturn}

% =====================================================================
\section*{\sffamily Ladder 9 \textbullet\ Periodic trends, explained Coulombically}
\ced{1.7}

Across a period, protons are added to the \emph{same} shell: nuclear
charge rises, shielding barely changes, so the atom contracts and holds
its electrons more tightly. Down a group, a new shell is added: distance
and shielding both rise, so the atom expands and holds its outer electrons
more loosely.

\begin{workedexample}[I do: ranking and explaining]
Rank \ce{Na}, \ce{Mg} and \ce{K} by first ionization energy, and justify.

\textbf{Order: \ce{K} $<$ \ce{Na} $<$ \ce{Mg}} (419, 496, 738
\si{\kilo\joule\per\mole}).

\textbf{\ce{Na} against \ce{Mg}} --- same shell, but magnesium has one
more proton. The extra nuclear charge holds the valence electrons more
tightly, so more energy is needed: \ce{Mg} is higher.

\textbf{\ce{Na} against \ce{K}} --- potassium's valence electron is in the
next shell out, further from the nucleus and shielded by an extra filled
shell. The attraction is weaker, so it is easier to remove: \ce{K} is
lower.

\textbf{What earns zero:} ``\ce{Mg} wants to lose two electrons'' or
``\ce{K} is more reactive''. Name charge, distance, shielding --- nothing
else counts.
\end{workedexample}

\begin{yourturn}[YOUR TURN 9 --- four questions]
\begin{enumerate}[leftmargin=*,itemsep=0.5em]
  \item Rank \ce{Li}, \ce{Na}, \ce{K} by atomic radius, smallest first.
        \workspace[1.4cm]
  \item Rank \ce{Na}, \ce{Mg}, \ce{Al} by atomic radius, smallest first.
        \workspace[1.4cm]
  \item The first ionization energy of sulfur (1000) is \emph{lower} than
        that of phosphorus (1012), breaking the trend. Explain.
        \workspace[2.4cm]
  \item Why is the \emph{second} ionization energy of sodium enormously
        larger than its first? \workspace[2.2cm]
\end{enumerate}
\selfcheck{(a) \ce{Li} $<$ \ce{Na} $<$ \ce{K} \quad
(b) \ce{Al} $<$ \ce{Mg} $<$ \ce{Na} \quad (c) the paired \ce{3p}
electron is repelled \quad (d) the second comes from a core shell}
\keyonly{\begin{apcanswer}(a) \ce{Li} $<$ \ce{Na} $<$ \ce{K}. Each step
down adds a shell, so distance and shielding both increase.
(b) \ce{Al} $<$ \ce{Mg} $<$ \ce{Na}. Across the period the nuclear charge
rises while the shell stays the same, pulling the cloud inward.
(c) Phosphorus is \ce{3p^3}, with one electron in each of the three
\ce{3p} orbitals. Sulfur is \ce{3p^4}, so one orbital holds a
\emph{pair}. Two electrons confined to the same orbital repel each other
strongly, raising that electron's energy and making it easier to remove ---
enough to outweigh sulfur's extra proton. (This is an argument about
electron--electron repulsion, \emph{not} about ``half-filled subshells
being stable''.)
(d) The first electron removed is sodium's single \ce{3s} valence
electron, far out and well shielded. Once it is gone, \ce{Na+} has a
filled \ce{2p} shell, and the next electron must come from that
\textbf{core} shell --- much closer to the nucleus and shielded by only two
electrons. Far greater effective nuclear charge at a far shorter distance
means an enormously larger energy.\end{apcanswer}}
\end{yourturn}

% =====================================================================
\section*{\sffamily Ladder 10 \textbullet\ Valence electrons and ionic compounds}
\ced{1.8}

A main-group element's group number gives its valence-electron count, and
that fixes the ion it forms. Combine ions so the charges cancel; the
lattice energy that results is Coulombic.

\begin{workedexample}[I do: from groups to a formula and a comparison]
Magnesium reacts with nitrogen. Give the ions, the formula, and predict
whether the compound's lattice energy exceeds that of \ce{NaF}.

\textbf{Ions:} Mg is group 2, two valence electrons, forms \ce{Mg^2+}.
N is group 15, five valence electrons, forms \ce{N^3-}.

\textbf{Formula:} the charges must cancel. Three \ce{Mg^2+} ($+6$) with
two \ce{N^3-} ($-6$) gives \textbf{\ce{Mg3N2}}.

\textbf{Lattice energy against \ce{NaF}:} compare Coulombically. The
attraction is proportional to the product of the ionic charges. For
\ce{Mg3N2} that product is $(2)(3) = 6$; for \ce{NaF} it is
$(1)(1) = 1$. \ce{Mg3N2} therefore has much the larger lattice energy.
\end{workedexample}

\begin{yourturn}[YOUR TURN 10 --- four questions]
\begin{enumerate}[leftmargin=*,itemsep=0.5em]
  \item Formula of the compound of aluminium and oxygen:
        \workspace[1.4cm]
  \item Formula of the compound of potassium and sulfur:
        \workspace[1.4cm]
  \item Which has the larger lattice energy, \ce{MgO} or \ce{NaF}?
        Justify Coulombically. \workspace[2.4cm]
  \item \ce{MgO} melts at \SI{2852}{\celsius}, \ce{NaF} at
        \SI{993}{\celsius}. Explain how this follows from your answer to
        (c). \workspace[2.0cm]
\end{enumerate}
\selfcheck{(a) \ce{Al2O3} \quad (b) \ce{K2S} \quad (c) \ce{MgO} --- charge
product 4 vs 1, and shorter distance \quad (d) stronger attractions take
more energy to overcome}
\keyonly{\begin{apcanswer}(a) \ce{Al^3+} and \ce{O^2-}:
$2(+3) + 3(-2) = 0$, so \textbf{\ce{Al2O3}}.
(b) \ce{K+} and \ce{S^2-}: $2(+1) + (-2) = 0$, so \textbf{\ce{K2S}}.
(c) \textbf{\ce{MgO}}, on both Coulombic counts. \textbf{Charge:} the
product is $(+2)(-2) = 4$ against \ce{NaF}'s $(+1)(-1) = 1$.
\textbf{Distance:} \ce{Mg^2+} and \ce{O^2-} are the smaller pair, so the
internuclear separation is shorter ($\approx$\SI{212}{\pico\metre}
against $\approx$\SI{235}{\pico\metre}) and the attraction, going as
$1/d$, is stronger again.
(d) Melting an ionic solid means giving the ions enough energy to break
free of the lattice's electrostatic attractions. \ce{MgO} has much the
stronger attractions, so far more energy --- a far higher temperature ---
is required.\end{apcanswer}}
\end{yourturn}

% =====================================================================
\section*{\sffamily Where you stand}

Tick a ladder only if all four YOUR TURN questions were right on the first
attempt.

\renewcommand{\arraystretch}{1.35}
% Ragged-right on the last column: justifying a 4.8cm measure with short
% phrases produces underfull hboxes for no visual gain.
\noindent\begin{tabular}{@{}c p{6.2cm} c >{\raggedright\arraybackslash}p{4.8cm}@{}}
\toprule
\textbf{First try?} & \textbf{Skill} & \textbf{Ladder} &
  \textbf{If not, re-read\ldots}\\
\midrule
$\square$ & moles, mass, particles & 1 & divide by $M$, then $\times N_A$ \\
$\square$ & molar mass with brackets & 2 & the subscript multiplies all \\
$\square$ & average atomic mass & 3 & weight by abundance \\
$\square$ & percent composition & 4 & per mole; must sum to 100 \\
$\square$ & empirical \& molecular & 5 & $\div$ smallest, then $\times$ \\
$\square$ & mixtures & 6 & no formula; separable by property \\
$\square$ & electron configurations & 7 & adjust for charge first \\
$\square$ & PES spectra & 8 & heights sum to $Z$ \\
$\square$ & periodic trends & 9 & charge, distance, shielding \\
$\square$ & ionic compounds & 10 & cancel charges; Coulomb's law \\
\bottomrule
\end{tabular}

\begin{apcnote}[Scoring yourself honestly]
10/10: move on to the unit worksheets and the free-response set.
7--9: solid --- redo the missed ladders tomorrow, not today, and from a
blank page. 6 or fewer: the recurring root causes in this unit are exactly
three --- (1) comparing \emph{masses} where the question needs
\emph{moles}, (2) forgetting that a subscript outside a bracket multiplies
everything inside it, and (3) explaining a periodic trend by what an atom
``wants'' instead of by nuclear charge, distance and shielding. Fix those
three and most of these ladders fall together.
\end{apcnote}

\end{document}
